
\lesson{11/11}



\[
    Z_0[J] = \sum_{n=0}^{\infty} \frac{i^N}{N!}\langle 
    G(1,2...N)J_1...J_N    \rangle_{1...N}
\]
And so the Green function, as seen before, is
\[
    G(1...N) = \eval{(-i)^N\frac{\delta^N}{\delta J_1...\delta J_N}Z_0[J]
}_{J_i=0}\]
\textbf{Explicit examples}\\
\[
    G_1(x_1)=\eval{-i\frac{\delta }{\delta J_1} Z_0[J] }_{J=0}= \eval{(-i)\left(-
    \frac{1}{2}\right)\frac{\delta }{\delta J_1}\left(\langle J(x)\Delta (x-y)
    J(y)\rangle_{xy}\right)Z_0[J]
}_{J=0}\]
the functional derivative acts on the exponent, yielding
\[
    G_1(x_1) = \eval{i \langle \underbrace{\frac{\delta J(x)}{\delta J(x_1)}}_{=
    \delta^4(x-x_1)}\Delta (x-y)J(y)\rangle_{xy}Z_0[J]}_{J=0}=
    \eval{i\langle \Delta (x_1-y)J(y)\rangle_yZ_0[J]}_{J=0}=0
\]

If there are two fields
\begin{gather*}
    G_2(x_1,x_2) = \eval{-\frac{\delta^2}{\delta J_1\delta J_2}}_{J=0}=
    \eval{\frac{\delta }{\delta J_2}\left(\langle \Delta (x-y)J(y)\rangle_y Z_0[J]
    \right)}_{J=0}=\\=
    \eval{-\Delta (x_1-x_2 )Z_0[J] - \langle \Delta (x_1-y)J(y)\rangle_y\underbrace{
    \frac{\delta }{\delta J_2}Z_0[J]}_\text{prev. res.}}_{J=0}=\\=
    \eval{-\Delta (x_1-x_2 )Z_0[J] + \langle \Delta (x_1-y)J(y)\rangle_y
    +\langle \Delta (x_2-w)J(w)\rangle_w Z_0[J]}_{J=0}=\\=-\Delta (x_1-x_2)
\end{gather*}
The $G_3$ propagator is again null, since an odd number of derivatives brings down
$J$ terms that vanish when setting $J=0$ (much like the Gaussian integrals).

For $4$ fields,
\begin{gather*}
    \eval{G_4 = \frac{\delta ^4}{\delta J_1\delta J_2\delta J_3}Z_0[J]}_{J=0}=
    \eval{\frac{\delta ^3}{\delta J_1\delta J_2\delta J_3}\left(
    \langle -\Delta (x_4-y)J(y)\rangle_y Z_0[J]
    \right)}_{J=0}=\\
    =\frac{\delta^2}{\delta J_1\delta J_2}\eval{\left(
    \left(-\Delta_{4,3}+\langle \Delta_{4,y}J_y\rangle_y\langle \Delta_{3w}J(w)
    \rangle_W\right)Z_0[J]\right)}_{J=0} = \\
    =\eval{\frac{\delta }{\delta J_1}\left(\left(
    \Delta_{43}\langle \Delta_{2y}J(y)\rangle_y + \Delta_{42}\langle \Delta_{3w}
    J(w)\rangle_w + \langle \Delta_{4y}J(y)\rangle_y\Delta_{32}-
    \langle\Delta_{4y}J(y)\rangle_y\langle \Delta_{3w}J(w)\rangle_w\langle 
    \Delta_{2z}J(z)\rangle_z\right)Z_0[J]\right)}_{J=0}=\\
    =\footnote{mi son perso...}
    \Delta_{32}\Delta _{41}-\Delta _{32}\langle \Delta _{4y}J_y\rangle_y
    \langle \Delta_{1w}J_w\rangle_w [... termini]
\end{gather*}
these are all possible derivatives of $4$-th order, by setting $J=0$, all the
$\langle ...J\rangle_{i}\Delta_{kj}$ will vanish, so
\[
G_4 = \Delta_{43}\Delta_{21} + \Delta_{42}\Delta_{31}+\Delta_{41}\Delta_{32}
\]
A trick to take these derivatives is to only compute those that will result in
terms independent of $J$!.

The result is known! These are the contractions seen in IQFT!
\[
    G_4 = \bra{0}T\left\{\phi _1\phi _2\phi _3\phi _4\right\}\ket{0} = 
    \bra{0} N\left\{\phi _1...\phi _4 + \text{contractions}\right\}\ket{0}=
    \bra{0}\overline{\phi_1\phi _2}\,\overline{\phi_3\phi _4} + ...\ket{0}
\]
and the result is the same!! This means that Wick's theorem is built-in for the
path integral formulation.

Is there a formal way to exclude non-connected Feynman Diagrams?

Consider
\[
    Z_0[J]= e^{-W_0[J]}
\]
where $W_0[J]$ is called the \textit{Free energy operator}; so
\[
    W_0[J] = \frac{1}{2}\int_{}^{} \dd ^4x\dd^4y J(x)\Delta(x-y)J(y)
\]
This generates \textbf{only connected Feynman diagrams}:
\[
W_0[J]=\sum_{N=0}^{\infty} \frac{1}{N!}\langle J_1...J_NG^C(1...N)\rangle_{1..N}
\implies \eval{G^C(1...N)=\frac{\delta^N}{\delta J_1...\delta J_N}W_0[J]}_{J=0}
\]
for the free theory this is due to the fact that $W_0$ is quadratic in $J$ and
all higher or lower order derivatives are null.

\section{Interacting theories: introduction}
The interacting action is
\[
    S=\int_{}^{} \dd ^4x\left(\frac{1}{2}\pd_\mu\phi \pd^\mu\phi -\frac{m^2}{2}
    \phi^2 -V(\phi)\right) = S_0 - \int_{}^{} \dd ^4x V(\phi)
\]
The partition function for interacting theories is
\begin{gather*}
    Z[J]=\frac{\int \mathcal{D}\phi
    \exp\left[iS+i\langle J\phi \rangle_x\right]}
    {\int \mathcal{D}\phi e^{iS_0}} = 
    \frac{\int \mathcal{D}\phi
    \exp\left[ -i\int_{}^{} \dd^4x V(\phi) \right]
    \exp\left[iS_0+i\langle J\phi \rangle_x\right]}
    {\int \mathcal{D}\phi e^{iS_0}} = \\
    =\mathcal{N} \int \mathcal{D}\phi e^{-i\langle V\rangle}e^{iS_0+i\langle 
    \phi J\rangle}
\end{gather*}
considering
\[
    e^{-i\langle V\rangle }=\sum_{n}^{} c_n\phi^n 
\]
this becomes, much like the gaussian integrals,
\begin{gather*}
    Z[J]=\mathcal{N}\sum_{n}^{} c_n \int \mathcal{D}\phi \phi^ne^{iS_0+i\langle
    \phi J\rangle } = \mathcal{N}\sum_{n}^{} c_n\int \mathcal{D}\phi \left(
    -i\frac{\delta}{\delta J}\right)^ne^{iS_0+\langle \phi J\rangle }=\\
    =\exp\left[ -i\langle V\left(-i\frac{\delta }{\delta J}\right)
    \rangle_x\right]\underbrace{\frac{\int \mathcal{D}\phi \exp\left[ iS_0
    +i\langle J\phi\rangle_x  \right]}{\int \mathcal{D}\phi e^{iS_0}}}_{:=Z_0[J]}=
    \\=\exp\left[ -i\langle V\left(-i\frac{\delta }{\delta J}\right)
    \rangle_x\right] Z_0[J]
\end{gather*}
How can this be computed? \textbf{Always with perturbation theory}, by
setting 
\[
    V(\phi)\to\Lambda V(\phi),\,\Lambda <<1
\]
and hoping the physical predictions are independent of $\Lambda $. 

Like in free theories, the free energy $W[J]$ is introduced: letting the
exponential operator be $\delta $ the partition function is
\begin{gather*}    
    Z[J]=e^{-W[J]} = \delta e^{-W_0} =
    e^{-W_0}e^{W_0}(1+\delta-1) e^{-W_0}
\end{gather*}
taking the $-\ln$ yields
\begin{gather*}
    W[J] = -\ln Z[J] = W_0-\ln\left(e^{W_0}(1+(\delta-1))e^{-W_0}\right)=
    W_0-\ln\left(1+\underbrace{e^{W_0}(\delta -1)e^{-W_0}}_{:=A}\right)
\end{gather*}
Since $\Lambda <<1$,
\[
    \delta =e^{-i\Lambda \langle V\rangle } = 1+O(\Lambda)
\]
and $A$ is small, so the $\ln$ can be expanded in Taylor
\[
    W[J]=W_0 - A+\frac{A^2}{2}+\frac{A^3}{3}...
\]
since $A$ depends on $\delta $, it can be expanded in powers of $\Lambda $:
\[
    A=\alpha _1\Lambda +\alpha _2\Lambda ^2+...
\]
and the free energy turns out to be
\begin{gather*}
    W[J]=W_0 - \Lambda \alpha _1-\Lambda ^2\alpha _2-\Lambda ^3a_3+
    \frac{1}{2}\left(\Lambda \alpha _1+\Lambda^2\alpha _2+\Lambda ^3\alpha _3
    \right)^2+
    \frac{1}{3}\left(\Lambda \alpha _1+\Lambda^2\alpha _2+\Lambda ^3\alpha _3
    \right)^3+...
\end{gather*}
reordering and truncating at $\Lambda ^3 $ yields
\[
    W[J]=W_0[J]-\Lambda \alpha _1-\Lambda ^2\left(\alpha _2-\frac{\alpha _1^2}{2}
    \right)-\Lambda ^3\left(\alpha _3 - \alpha _1\alpha _2+\frac{\alpha_1^3}{3}
    \right) + O(\Lambda^4)
\]
How can these coefficients be computed? Recalling
\[
    A=e^{W_0}(1-\delta)e^{-W_0}
\]
the first term is\footnote{do the calculation yourself}
\[
    \alpha _1 = -ie^{W_0}\left\langle{V\left(-i\frac{\delta }{\delta J}\right)}
    \right\rangle e^{-W_0}
\]
\section{$\phi^4$ interaction}

\textbf{If} the potential is
\[
    V(\phi) = \frac{\Lambda }{4!}\phi ^4
\]
the operator $\delta$ has the form
\[
    \delta = e^{-i\Lambda/4! \left\langle \frac{\delta^4}{\delta J^4}
    \right\rangle_x }=1-\frac{i\Lambda}{4!} \left\langle \frac{\delta^4}
    {\delta J ^4}\right\rangle_x
    -\frac{1}{2}\frac{\Lambda ^2}{(4!)^2}\left\langle \frac{\delta^4}
    {\delta J ^4}\right\rangle_x\left\langle \frac{\delta^4}{\delta J ^4}
    \right\rangle_y+O(\Lambda^3)
\]
